\section{Reductions at syntactic bidegree two}\label{sec:degree-two}

This section gives replacements that preserve shared wires and do not
increase the number of gates in two specified classes of circuits.
Let $G=(X\sqcup Y,E)$ be bipartite.  A positive polynomial is
\emph{safe} if every one of its monomials has support containing an
edge of $G$. Syntactic bidegree bounds the formal degree separately in
$X$ and $Y$: variable inputs have degrees $(1,0)$ or $(0,1)$, a product
adds the two input pairs, and an addition takes their componentwise
maximum. A useful gate with an addition-only route from that gate
to the output of a circuit for $f_G$ is safe.  In a variable-input
circuit of syntactic bidegree at most $(2,2)$, every useful gate of
bidegree $(2,2)$ has only additions above it.

\begin{lemma}[Four existing pure-side clauses]
Suppose a useful safe gate of a monotone circuit computing $f_G$
has Boolean value
\[
 g=A(X)B(X)C(Y)D(Y),
\]
where all four factors are existing nonempty linear disjunction wires.
If $G$ is $C_4$-free, then $g$ can be replaced by zero or by a single
star product using an existing factor wire and one variable, without
increasing the circuit size or changing its output function $f_G$.
Sharing of all four factor wires with the rest of the circuit is
permitted.
\end{lemma}
\begin{proof}
Write $A,B,C,D$ also for their variable supports, and put
$I=A\cap B$, $J=C\cap D$.  The gate accepts the edge assignment
$\{x,y\}$ exactly when $(x,y)\in I\times J$.  If either intersection
is empty, $g$ is false on every edge assignment. Replacing it by zero
leaves all these assignments unchanged and cannot create a false
positive, so the output remains $f_G$.  Otherwise safety gives $I\times J\subseteq E$.
The absence of a four-cycle implies $|I|=1$ or $|J|=1$.

Suppose $I=\{x\}$.  If both $C$ and $D$ contain a nonneighbor of $x$,
choose $y\in C\setminus N(x)$ and $z\in D\setminus N(x)$.
The gate accepts $\{x,y,z\}$, an independent set, which is impossible.
Hence $C\subseteq N(x)$ or $D\subseteq N(x)$.
Replace $g$ by $xC$ or $xD$, respectively, reusing the gate and the
specified factor wire.  This is a safe star and accepts every old
edge witness because $J\subseteq C,D$.  The case $|J|=1$ is symmetric.
The replacement is safe and still accepts every edge assignment
accepted by $g$. Every new output monomial passing through the replacement
therefore contains an edge, while all other output monomials are
unchanged. On each edge assignment the replacement cannot change a true
gate value to false, so monotonicity preserves every output edge witness.
Thus the new output is safe and still accepts every edge, hence equals
$f_G$. The pure-side linear factor wires cannot depend on $g$, whose
value requires variables of both sides, so this rewiring creates no
cycle. Zero propagation and pruning add no gates.
\end{proof}

The replacement uses the existing factor wires.  Merely expanding a sum
of bilinear products into many four-clause products would not preserve
the gate count.  The following simultaneous replacement avoids such
an expansion for a larger, explicitly restricted architecture.

\begin{theorem}[A shared bilinear two-level class]
Suppose a binary monotone circuit has the following form:
\begin{enumerate}
\item Its first product gates are $R_i=A_i(X)B_i(Y)$, with pure-side
linear disjunction operands, shared arbitrarily.
\item Its second product gates multiply two disjunctions of the $R_i$.
\item Its output is a disjunction of first and second product gates.
\end{enumerate}
All disjunction networks may have arbitrary fanout and sharing.
Suppose the circuit computes $f_G$ for a nonempty $C_4$-free bipartite
graph.  For every unsafe rectangle put
\[
 U_i=\{x\in A_i:B_i\subseteq N(x)\}.
\]
If $|U_i|\leq1$ for every unsafe rectangle, the circuit can be
replaced without growth by a quadratic circuit with edge support $E$.
In particular, $|B_i|\geq2$ for every unsafe rectangle is sufficient.  Consequently its size is at
least $|E|+\tau(G)-1$, using the established exact quadratic theorem (Theorem~\ref{thm:granted}).
The symmetric assertion holds with $X$ and $Y$ exchanged.
\end{theorem}
\begin{proof}
Keep every safe rectangle.  For an unsafe rectangle $R_i$, the
hypothesis gives at most one $x\in A_i$ with $B_i\subseteq N(x)$.
If that $x$ exists, replace $R_i$ by $xB_i$, using its existing
product gate and its existing $B_i$ wire.  If $U_i$ is empty, replace
it by zero.  Replace every second product gate by a disjunction gate.
All remaining first products are safe, so the resulting output is
safe; the replacements introduce no gate and do not change any shared
linear wire.

We show that every old edge witness remains.  A first product with
an addition-only route to the output was safe, and was kept.  If an
edge $xy$ was witnessed at a second product, choose first rectangles
$R_i=A_iB_i$ and $R_j=A_jB_j$ in its two operands which both accept
$\{x,y\}$.  Their product is safe.  If either rectangle is safe,
it was kept and still witnesses the edge at the new disjunction.
Otherwise both rectangles are unsafe.
Safety on assignments with sole $X$-variable $x$ implies
\[
 B_i\subseteq N(x)\quad\text{or}\quad B_j\subseteq N(x):
\]
if neither inclusion held, choosing one nonneighbor from each
$Y$-support would make their product true on an independent set.
In the first case the unique selected centre of $R_i$ is $x$;
in the second the unique selected centre of $R_j$ is $x$.
Thus the new disjunction accepts $xy$.  Every edge is preserved,
and safety proves equality with $f_G$.

Propagate the introduced zeros, bypass OR gates with a zero operand,
and remove unused gates; no new constant gate need remain.
The circuit is an OR of bilinear star products, hence
its arithmetic polynomial is quadratic with exactly the required
edge support.  Apply the exact quadratic theorem, cited in Theorem~\ref{thm:granted}.
Finally, in a $C_4$-free graph two universal rows would force
$|B_i|\leq1$, proving the stated sufficient condition.
\end{proof}

In particular, within this architecture a counterexample must contain
an unsafe rectangle $A\times\{y\}$ with at least two vertices of
$A$ adjacent to $y$, and, symmetrically, an unsafe rectangle
$\{x\}\times B$ with at least two vertices of $B$ adjacent to $x$.
Each of these rectangles also contains a nonneighbor, by unsafety.
They are necessarily distinct.
This theorem does not include mixed-degree operands or all circuits
with two levels of AND gates.

\begin{lemma}[Disjoint bad endpoint sets]
Let $P,Q$ be positive bilinear polynomials and suppose $PQ$ is safe.
The nonedge pairs occurring in $P$ and those occurring in $Q$ have
disjoint endpoint sets on each side of the bipartition.
\end{lemma}
\begin{proof}
If $xy$ occurs in $P$, $xz$ occurs in $Q$, and both pairs are
nonedges, the product contains $x^2yz$, whose support is independent.
This contradicts safety.  The argument on the other side is the same.
No four-cycle assumption is needed.
\end{proof}

\begin{example}[Both factors can be unsafe]
Let $G$ be the matching with edges $x_1y_1,x_2y_2$, and put
\[
 P=x_1y_1+x_1y_2+x_2y_2,\qquad
 Q=x_1y_1+x_2y_1+x_2y_2.
\]
Both factors contain a nonedge monomial, but their product is safe
and witnesses both matching edges.  Indeed,
\[
\begin{split}
PQ={}&x_1^2y_1^2+x_1^2y_1y_2+x_1x_2y_1^2
 +3x_1x_2y_1y_2\\
 &+x_1x_2y_2^2+x_2^2y_1y_2+x_2^2y_2^2.
\end{split}
\]
Every displayed term contains one of the matching edges.
Thus safety of a bilinear product does not imply safety of an
operand.  This is a counterexample to that proof shortcut, not to
the circuit lower bound.
\end{example}

